% GENERATED FILE. Edit canonical lessons, then run npm run generate:books.
% canonical-source-hash: fnv1a64:45e09e97be0ffc5f
% canonical-lessons: TA-C443-oti, TA-C443-norunku, TA-C443-virical, TA-C443-ottai, TA-C443-kiru

\chapter{Cracks and Holes}
\label{ch:ta-cracks-and-holes}
\hlchaptermodality{443}

\begin{chapteropening}
\textbf{By the end of this chapter:} \emph{I can talk about things breaking off and shattering, cracks, holes and scratches.}

Complete the last lesson of chapter 443: I can talk about things breaking off and shattering, cracks, holes and scratches.
\end{chapteropening}

\section[\texorpdfstring{\ta{ஒடி} (oṭi)}{ஒடி (oṭi)}]{\ta{ஒடி} (oṭi) — to break off}

\label{lesson:TA-C443-oti}



\begin{quote}
Before the new one: say the Tamil for to surround, then the Tamil for to swing about.
\end{quote}

\subsection*{You'll want to know: \ta{ஒடி}}
\textbf{\ta{ஒடி}} — \emph{oṭi} — \textquotedblleft{}to break off\textquotedblright{}.

\begin{quote}
\ta{புயலில்} \ta{மரக்கிளை} \ta{ஒடிந்தது}.
\end{quote}

\begin{quote}
\emph{puyalil marakkiḷai oṭintatu.}
\end{quote}

\begin{quote}
A branch broke off in the storm.
\end{quote}

Past (it …): \ta{ஒடிந்தது}, \emph{oṭintatu}.

\ta{ஒடிந்தது} is it broke off; \ta{ஒடித்தேன்}, I broke it off.

\begin{cousinweb}
\textbf{\ta{ஒடி}} (\emph{oṭi}) is native Dravidian (DEDR 946), \textquotedblleft{}to break, to break off, to snap\textquotedblright{}.
\end{cousinweb}

\begin{grammarlens}[title={What you've built}]
One more word for getting around.
\end{grammarlens}

\subsection*{Your turn}
\begin{itemize}
\raggedright
  \item \emph{Say it:} \emph{oṭi}
  \item \emph{Say it:} \emph{oṭi}, once more
  \item \emph{Say it:} \emph{puyalil marakkiḷai oṭintatu.} — A branch broke off in the storm.
  \item \emph{Recall:} say the Tamil for to surround, then the Tamil for to swing about, then say \emph{oṭi} again
\end{itemize}

\begin{tcolorbox}[breakable,colback=teal!4,colframe=teal!35!black,title={Before you move on}]
What does \ta{ஒடி} mean? (\textquotedblleft{}To break off\textquotedblright{}.) Say it once more.
\end{tcolorbox}
\section[\texorpdfstring{\ta{நொறுங்கு} (noṟuṅku)}{நொறுங்கு (noṟuṅku)}]{\ta{நொறுங்கு} (noṟuṅku) — to shatter; to be smashed}

\label{lesson:TA-C443-norunku}



\begin{quote}
Before the new one: say the Tamil for to swing about, then the Tamil for to break off.
\end{quote}

\subsection*{You'll want to know: \ta{நொறுங்கு}}
\textbf{\ta{நொறுங்கு}} — \emph{noṟuṅku} — \textquotedblleft{}to shatter; to be smashed\textquotedblright{}.

\begin{quote}
\ta{கண்ணாடி} \ta{கீழே} \ta{விழுந்து} \ta{நொறுங்கியது}.
\end{quote}

\begin{quote}
\emph{kaṇṇāṭi kīḻē viḻuntu noṟuṅkiyatu.}
\end{quote}

\begin{quote}
The glass fell down and shattered.
\end{quote}

Past (it …): \ta{நொறுங்கியது}, \emph{noṟuṅkiyatu}.

\ta{நொறுக்கு} is to smash something.

\begin{cousinweb}
\textbf{\ta{நொறுங்கு}} (\emph{noṟuṅku}) is native Dravidian (DEDR 3728), \textquotedblleft{}to be broken, crushed, smashed\textquotedblright{}.
\end{cousinweb}

\begin{grammarlens}[title={What you've built}]
One more word for getting around.
\end{grammarlens}

\subsection*{Your turn}
\begin{itemize}
\raggedright
  \item \emph{Say it:} \emph{noṟuṅku}
  \item \emph{Say it:} \emph{noṟuṅku}, once more
  \item \emph{Say it:} \emph{kaṇṇāṭi kīḻē viḻuntu noṟuṅkiyatu.} — The glass fell down and shattered.
  \item \emph{Recall:} say the Tamil for to swing about, then the Tamil for to break off, then say \emph{noṟuṅku} again
\end{itemize}

\begin{tcolorbox}[breakable,colback=teal!4,colframe=teal!35!black,title={Before you move on}]
What does \ta{நொறுங்கு} mean? (\textquotedblleft{}To shatter; to be smashed\textquotedblright{}.) Say it once more.
\end{tcolorbox}
\section[\texorpdfstring{\ta{விரிசல்} (virical)}{விரிசல் (virical)}]{\ta{விரிசல்} (virical) — a crack}

\label{lesson:TA-C443-virical}



\begin{quote}
Before the new one: say the Tamil for to break off, then the Tamil for to shatter.
\end{quote}

\subsection*{You'll want to know: \ta{விரிசல்}}
\textbf{\ta{விரிசல்}} — \emph{virical} — \textquotedblleft{}a crack\textquotedblright{}.

\begin{quote}
\ta{சுவரில்} \ta{ஒரு} \ta{விரிசல்} \ta{தெரிகிறது}.
\end{quote}

\begin{quote}
\emph{cuvaril oru virical terikiṟatu.}
\end{quote}

\begin{quote}
A crack can be seen in the wall.
\end{quote}

Built on \ta{விரி}, to split open; \ta{உறவில்} \ta{விரிசல்} is a rift between people.

\begin{cousinweb}
\textbf{\ta{விரிசல்}} (\emph{virical}) is native Dravidian (DEDR 5411), \textquotedblleft{}a split, a crack\textquotedblright{}, the noun of \ta{விரி}, \textquotedblleft{}to open, to split\textquotedblright{}.
\end{cousinweb}

\begin{grammarlens}[title={What you've built}]
One more word for getting around.
\end{grammarlens}

\subsection*{Your turn}
\begin{itemize}
\raggedright
  \item \emph{Say it:} \emph{virical}
  \item \emph{Say it:} \emph{virical}, once more
  \item \emph{Say it:} \emph{cuvaril oru virical terikiṟatu.} — A crack can be seen in the wall.
  \item \emph{Recall:} say the Tamil for to break off, then the Tamil for to shatter, then say \emph{virical} again
\end{itemize}

\begin{tcolorbox}[breakable,colback=teal!4,colframe=teal!35!black,title={Before you move on}]
What does \ta{விரிசல்} mean? (\textquotedblleft{}A crack\textquotedblright{}.) Say it once more.
\end{tcolorbox}
\section[\texorpdfstring{\ta{ஓட்டை} (ōṭṭai)}{ஓட்டை (ōṭṭai)}]{\ta{ஓட்டை} (ōṭṭai) — a hole; a leak}

\label{lesson:TA-C443-ottai}



\begin{quote}
Before the new one: say the Tamil for to shatter, then the Tamil for a crack.
\end{quote}

\subsection*{You'll want to know: \ta{ஓட்டை}}
\textbf{\ta{ஓட்டை}} — \emph{ōṭṭai} — \textquotedblleft{}a hole; a leak\textquotedblright{}.

\begin{quote}
\ta{படகில்} \ta{ஓட்டை} \ta{இருந்ததால்} \ta{தண்ணீர்} \ta{உள்ளே} \ta{வந்தது}.
\end{quote}

\begin{quote}
\emph{paṭakil ōṭṭai iruntatāl taṇṇīr uḷḷē vantatu.}
\end{quote}

\begin{quote}
Water came in because there was a hole in the boat.
\end{quote}

\ta{ஓட்டை} \ta{விழு} is to get a hole.

\begin{cousinweb}
\textbf{\ta{ஓட்டை}} (\emph{ōṭṭai}) is native Dravidian (DEDR 946), \textquotedblleft{}a hole, a crack in a vessel, a leak\textquotedblright{}.
\end{cousinweb}

\begin{grammarlens}[title={What you've built}]
One more word for getting around.
\end{grammarlens}

\subsection*{Your turn}
\begin{itemize}
\raggedright
  \item \emph{Say it:} \emph{ōṭṭai}
  \item \emph{Say it:} \emph{ōṭṭai}, once more
  \item \emph{Say it:} \emph{paṭakil ōṭṭai iruntatāl taṇṇīr uḷḷē vantatu.} — Water came in because there was a hole in the boat.
  \item \emph{Recall:} say the Tamil for to shatter, then the Tamil for a crack, then say \emph{ōṭṭai} again
\end{itemize}

\begin{tcolorbox}[breakable,colback=teal!4,colframe=teal!35!black,title={Before you move on}]
What does \ta{ஓட்டை} mean? (\textquotedblleft{}A hole; a leak\textquotedblright{}.) Say it once more.
\end{tcolorbox}
\section[\texorpdfstring{\ta{கீறு} (kīṟu)}{கீறு (kīṟu)}]{\ta{கீறு} (kīṟu) — to scratch}

\label{lesson:TA-C443-kiru}



\begin{quote}
Before the new one: say the Tamil for a crack, then the Tamil for a hole.
\end{quote}

\subsection*{You'll want to know: \ta{கீறு}}
\textbf{\ta{கீறு}} — \emph{kīṟu} — \textquotedblleft{}to scratch\textquotedblright{}.

\begin{quote}
\ta{பூனை} \ta{என்} \ta{கையைக்} \ta{கீறியது}.
\end{quote}

\begin{quote}
\emph{pūṉai eṉ kaiyaik kīṟiyatu.}
\end{quote}

\begin{quote}
The cat scratched my hand.
\end{quote}

Past (I …): \ta{கீறினேன்}, \emph{kīṟiṉēṉ}.

Also to cut a line: \ta{கோடு} \ta{கீறு}, to scratch a line.

\begin{cousinweb}
\textbf{\ta{கீறு}} (\emph{kīṟu}) is native Dravidian (DEDR 1623), \textquotedblleft{}to scratch, to draw lines\textquotedblright{}.
\end{cousinweb}

\begin{grammarlens}[title={What you've built}]
One more word for getting around.
\end{grammarlens}

\subsection*{Your turn}
\begin{itemize}
\raggedright
  \item \emph{Say it:} \emph{kīṟu}
  \item \emph{Say it:} \emph{kīṟu}, once more
  \item \emph{Say it:} \emph{pūṉai eṉ kaiyaik kīṟiyatu.} — The cat scratched my hand.
  \item \emph{Recall:} say the Tamil for a crack, then the Tamil for a hole, then say \emph{kīṟu} again
\end{itemize}

\begin{tcolorbox}[breakable,colback=teal!4,colframe=teal!35!black,title={Before you move on}]
What does \ta{கீறு} mean? (\textquotedblleft{}To scratch\textquotedblright{}.) Say it once more.
\end{tcolorbox}
